\documentclass[conference]{IEEEtran}

\usepackage[utf8]{inputenc}
\usepackage{amsmath,amssymb}
\usepackage{booktabs}
\usepackage{cite}
\usepackage{graphicx}
\usepackage{textcomp}

\begin{document}

\title{Multi-Seed Evaluation of Established Deep Neural Decoders on Low-Cost Motor-Imagery EEG}

\author{
\IEEEauthorblockN{Miguel Isidro Baez, Sharon Quispe Carhuapoma, Omar Perez Nestares,
Valeria Rivera Angles}
\IEEEauthorblockA{Universidad Peruana Cayetano Heredia, Universidad San Ignacio de Loyola,
and Universidad de San Martin de Porres, Lima, Peru}
}

\maketitle

\begin{abstract}
This study evaluates five established deep neural decoders for binary motor-imagery-versus-rest
classification on the MI-OpenBCI low-cost EEG dataset. EEGNet, FBCNet, ShallowConvNet,
EEGConformer, and EEGInceptionMI are trained under fixed, auditable recipes and repeated with five
optimization seeds. Participant-calibrated, five-fold within-session, and leave-one-subject-out
protocols are reported separately. External validation uses the task-matched, multi-session
Zhou2020 and Tavakolan2017 datasets, with leave-one-session-out evaluation. Model comparisons use
participants as the statistical unit, two-sided paired Wilcoxon tests, Holm correction, confidence
intervals, and paired rank-biserial effect sizes. The study is offline and does not estimate a
causal effect of hardware cost or end-to-end BCI response time. Confirmatory numerical results will
be inserted only after the complete five-seed experiment grid passes the prespecified integrity
checks.
\end{abstract}

\begin{IEEEkeywords}
motor imagery, electroencephalography, low-cost EEG, deep learning, cross-session evaluation,
reproducibility
\end{IEEEkeywords}

\section{Introduction}

Motor-imagery brain-computer interfaces decode changes in sensorimotor EEG associated with
imagined movement. Most established decoders are evaluated on laboratory-grade recordings, while
public low-cost datasets contain fewer participants and may have different noise characteristics.
This study therefore asks a narrower question: how do five representative deep decoders behave on
MI-OpenBCI under participant-calibrated and cross-participant protocols, and are the observed
patterns consistent in task-matched research-grade datasets?

The contribution is empirical. No new neural architecture is proposed. The study adds a
predeclared multi-seed design, participant-level paired inference, task-matched cross-session
validation, leakage checks, sample accounting, and a reproducible model-inference latency profile.
Differences between datasets are reported separately because the acquisitions do not isolate
hardware cost.

\section{Methods}

\subsection{Datasets}

MI-OpenBCI is the low-cost primary dataset. It contains 10 available participants, 15 EEG channels,
and a right-hand motor-imagery-versus-rest task acquired at 125 Hz. The original subject identifiers
are retained; S01 and S11 are absent from the distributed dataset. Exact trials and class counts are
generated from the loaded arrays rather than entered manually.

External validation uses Zhou2020~\cite{zhou} and Tavakolan2017~\cite{tavakolan} through MOABB. Both contain the events
\texttt{right\_hand} and \texttt{rest} and have repeated sessions. Zhou2020 contains 20 participants
and seven sessions. Tavakolan2017 contains 12 participants and four sessions. BNCI2014\_001 is not
part of the confirmatory profile because it does not contain a rest class; a binary left-versus-
right reduction would answer a different task.

\subsection{Preprocessing}

Every fold applies the same preprocessing to all five models. Signals are resampled to 128 Hz,
filtered from 1 to 40 Hz, then filtered from 8 to 30 Hz and standardized per channel using mean and
standard deviation estimated only from the training partition. Primary analyses do not reject ICA
components.

An exploratory sensitivity analysis fits FastICA only on training data. Components with kurtosis
greater than 10 are candidates and at most the two largest-kurtosis components are removed. The
pipeline records convergence, iteration count, component kurtosis, candidates, and exclusions.
This rule is not described as ocular or muscular component identification.

\subsection{Decoders and training}

The five decoders are EEGNet, FBCNet, ShallowConvNet implemented as Braindecode's
ShallowFBCSPNet, EEGConformer, and EEGInceptionMI. All implementations come from Braindecode
1.5.2. FBCNet uses six internal subbands restricted to 8--30 Hz. Table~\ref{tab:recipes} gives the
frozen recipes. The held-out test partition is never used for hyperparameter or epoch selection.

\begin{table}[t]
\caption{Predeclared training recipes. Every model uses 300 fixed epochs, batch size 64,
cosine annealing, and no early stopping.}
\label{tab:recipes}
\centering
\begin{tabular}{lcc}
\toprule
Model & Optimizer & Learning rate \\
\midrule
EEGNet & AdamW & $6.25\times10^{-4}$ \\
FBCNet & AdamW & $6.25\times10^{-4}$ \\
ShallowConvNet & AdamW & $6.25\times10^{-4}$ \\
EEGConformer & Adam & $5\times10^{-5}$ \\
EEGInceptionMI & Adam & $1\times10^{-3}$ \\
\bottomrule
\end{tabular}
\end{table}

Each experiment is repeated with optimization seeds 0, 1, 2, 3, and 4. The split seed is fixed at
2026. FastICA also uses a fixed seed derived from the split and fold. Consequently, dispersion over
the five repetitions represents optimization variability rather than a mixture of changing data
partitions and optimization.

\subsection{Evaluation protocols}

MI-OpenBCI uses three protocols that answer different questions: a stratified 70/30 held-out split
within the only available session, stratified five-fold within-session evaluation, and
leave-one-subject-out evaluation. LOSO measures transfer to an unseen participant within the same
dataset and is not interpreted as population-level generalization.

Zhou2020 and Tavakolan2017 use a stratified 70/30 split within each session and leave-one-session-
out evaluation within each participant. MI-OpenBCI has one session and therefore cannot provide a
cross-session result.

The augmentation analysis is separate from the primary full-trial table. All five models receive
the same two-second input under three training conditions: one center crop, two non-overlapping
windows, or six overlapping windows. Complete trials are split before segmentation. Test data
produce one prediction per original trial, so windows are never treated as independent statistical
observations.

\subsection{Metrics and statistical analysis}

Accuracy is the primary endpoint and Cohen's $\kappa$ is secondary. Cohen's $\kappa$ measures
observed agreement corrected for agreement expected by chance. Other metrics are descriptive.

The five optimization seeds are averaged within participant before inference. Each family is
defined by dataset, task, protocol, input condition, ICA policy, and metric. The ten model pairs are
compared using two-sided paired Wilcoxon tests followed by Holm correction. Tables include the mean
and median paired difference, a 95\% Student-$t$ interval for the mean paired difference, and paired
rank-biserial correlation. No Friedman test or bootstrap interval is used.

The augmentation family contains the two predeclared comparisons, non-overlapping minus center and
overlapping minus center, for each of the five models. Inference is blocked if any expected cell,
seed, participant cohort, code fingerprint, training recipe, environment identity, or data identity
is inconsistent.

\subsection{Latency scope}

Latency is measured after warm-up using the same PyTorch backend, device, channel count, temporal
length, and floating-point type for all models. Batch-one median and 95th percentile measure
individual model-forward latency. Batch-64 values describe amortized throughput. The measurement
excludes acquisition, window accumulation, transmission, preprocessing, user interface, and
actuation. It is not end-to-end BCI response time.

\section{Results}

\noindent\fbox{\parbox{0.95\columnwidth}{\textbf{Draft status.} Confirmatory results have not been
inserted. Complete the expected-cell audit and the five-seed analysis before replacing this box.}}


\section{Discussion}

The final discussion will distinguish the highest observed mean from statistically supported
pairwise differences. It will report optimization variability separately from between-participant
dispersion and will not treat five seeds as additional biological observations. Dataset-specific
results will remain separate because Zhou2020 and Tavakolan2017 do not provide a matched causal
comparison of low-cost and research-grade hardware.

\section{Limitations}

MI-OpenBCI contains 10 participants from one device and one session. External datasets add
cross-session evidence but not a second low-cost acquisition or a controlled hardware comparison.
The study is offline and does not measure information-transfer rate, false activations, calibration
burden, user workload, energy, or clinical effectiveness. The ICA sensitivity is reproducible but
does not assign physiological artifact labels. Model-forward latency does not represent complete
system latency.

\section{Conclusion}

Conclusions about decoder ranking, augmentation, and cross-session robustness remain pending until
the complete multi-seed results pass the prespecified integrity checks. The final text will be
generated from those verified outputs and will not reuse the historical single-seed results.

\begin{thebibliography}{00}

\bibitem{moabb}
V. Jayaram and A. Barachant, ``MOABB: trustworthy algorithm benchmarking for BCIs,''
\emph{Journal of Neural Engineering}, vol. 15, no. 6, 2018.

\bibitem{miopenbci}
V. Peterson, C. Galv{\'a}n, H. Hern{\'a}ndez, M. P. Saavedra, and R. Spies,
``A motor imagery vs. rest dataset with low-cost consumer grade EEG hardware,''
\emph{Data in Brief}, vol. 42, 2022.

\bibitem{eegnet}
V. J. Lawhern \emph{et al.}, ``EEGNet: a compact convolutional network for EEG-based
brain-computer interfaces,'' \emph{Journal of Neural Engineering}, vol. 15, no. 5, 2018.

\bibitem{shallow}
R. T. Schirrmeister \emph{et al.}, ``Deep learning with convolutional neural networks for EEG
decoding and visualization,'' \emph{Human Brain Mapping}, vol. 38, no. 11, 2017.

\bibitem{conformer}
Y. Song, Q. Zheng, B. Liu, and X. Gao, ``EEG Conformer: convolutional transformer for EEG decoding
and visualization,'' \emph{IEEE Transactions on Neural Systems and Rehabilitation Engineering},
vol. 31, 2023.

\bibitem{holm}
S. Holm, ``A simple sequentially rejective multiple test procedure,''
\emph{Scandinavian Journal of Statistics}, vol. 6, no. 2, 1979.

\bibitem{zhou}
Q. Zhou \emph{et al.}, ``Relative power correlates with the decoding performance of motor imagery
both across time and subjects,'' \emph{Frontiers in Human Neuroscience}, vol. 15, 2021.

\bibitem{tavakolan}
M. Tavakolan, Z. Frehlick, X. Yong, and C. Menon, ``Classifying three imaginary states of the same
upper extremity using time-domain features,'' \emph{PLOS ONE}, vol. 12, no. 3, 2017.

\end{thebibliography}

\end{document}
